\begin{recipe}{Andijviestamp met spekjes}
\kicker[Hoofdgerecht,Vlees,Nederlands,Doordeweeks]{Hoofdgerecht · vlees · Nederlands · doordeweeks}
\index[register]{Andijviestamp met spekjes}
\index[register]{Andijvie}
\index[register]{Spekjes}
\index[register]{Stamppot}

\meta{35 minuten}

\lettrine{A}{ndijviestamp} is een winterse stamppot: rauwe andijvie wordt
door de hete puree geroerd tot hij net geslonken is. Spekjes, mosterd en
kerrie maken het af.

\blockrule{Ingrediënten}
\begin{ingredients}
  \ing{1 kg}{kruimige aardappelen, geschild en in stukken}\index[register]{Aardappelen}
  \ing{600 g}{andijvie (zak), schoongemaakt}\index[register]{Andijvie}
  \ing{250 g}{gezouten spekblokjes}\index[register]{Spekjes}
  \ing{1 el}{mosterd}
  \ing{1 tl}{kerrie}
  \ingb{scheutje melk}
  \ingb{klontje boter}
  \ingb{versgemalen peper en zout}
\end{ingredients}

\blockrule{Bereiding}
\tip{Roer de andijvie er pas vlak voor het serveren door. Hoe langer de
stamppot daarna warm blijft, hoe meer vocht de andijvie afgeeft.}
\begin{steps}
  \item Kook de aardappelen in gezouten water gaar, ongeveer 20 minuten.
  \item Bak ondertussen de spekblokjes in een droge koekenpan knapperig.
        Voeg de kerrie toe en bak nog een minuut mee.
  \item Giet de aardappelen af en laat goed uitstomen. Stamp ze met een
        scheut melk, een klontje boter en de mosterd tot een stevige puree.
  \item Roer de andijvie in delen door de hete puree, tot hij geslonken is.
  \item Roer het spekmengsel erdoor. Breng op smaak met peper en zout.
\end{steps}
\end{recipe}
